\chapter{Archived Approaches and Project Evolution}
\label{ch:appendix_archive}

This appendix records the main lines of work that were explored seriously during the project but did not become part of the final reported pipeline. I include them for two reasons. First, they account for a substantial part of the technical work carried out during the project, and omitting them would give an inaccurate impression of how the final direction was reached. Second, several of these branches were reasonable ideas on paper. They were not abandoned because they were obviously misguided, but because they turned out to be harder to make reliable, less competitive than hoped, or simply less suitable for a final-year thesis whose main contribution needed to be coherent, reproducible, and explainable end to end.

All of the code discussed here is preserved in the repository under \path{ARCHIVE}, and can also be browsed on GitHub at \url{https://github.com/RohanDevereux/Bridging/tree/main/ARCHIVE}. The goal of this appendix is not to present those branches as polished alternatives to the final pipeline. The goal is to explain what they were trying to do, why they were attractive at the time, what was actually implemented, and why the project eventually moved on.

\section{Correction learning around fast affinity baselines}

One early project idea was to avoid predicting binding affinity entirely from scratch. Instead, the model could start from a fast existing estimator and learn a correction term on top of it. This is an attractive strategy because it combines a hand-designed or physics-inspired baseline with a statistical model that may learn where the baseline fails. Two branches were explored in this spirit: a contact-based baseline using PRODIGY \citep{Vangone2015Contacts} and an endpoint free-energy baseline using MM/GBSA \citep{Genheden2015MMGBSA}.

\subsection{PRODIGY}

The PRODIGY-related code is preserved under \path{ARCHIVE/PRODIGY}, or on GitHub at \url{https://github.com/RohanDevereux/Bridging/tree/main/ARCHIVE/PRODIGY}. This branch contains dataset parsing, chain-group normalization, estimate prefetching, and a runner that calls an external PRODIGY executable on a protein complex and parses the predicted $\Delta G$. Conceptually, this branch was motivated by the observation that contact-based affinity models are simple, interpretable, and surprisingly competitive for some protein--protein benchmarks. That motivation is well grounded in the literature: Vangone and Bonvin showed that interfacial contact counts and related structural descriptors can correlate meaningfully with experimental binding affinity in protein--protein complexes \citep{Vangone2015Contacts}.

In practical terms, the branch was intended to answer a simple question: if a contact-based structural predictor already gives a rough estimate, can a second model learn the residual error more easily than predicting affinity from raw structural input alone? This is a sensible idea, especially for an undergraduate project, because it keeps one foot in an interpretable baseline while still allowing machine learning to add value.

However, the preserved outputs in the repository were not convincing enough to justify making PRODIGY the core of the thesis. The committed exploratory estimates under \path{src/bridging/generatedData/PRODIGY} are sparse and were not produced under the final confirmatory protocol, so they should be treated as indicative rather than definitive. Even so, they point in the same direction as the qualitative decision that was made at the time. On the small preserved \path{PPB_Affinity_diverse100} exploratory subset, the matched cases in \path{PPB_Affinity_diverse100_prodigy_estimates.csv} give a Pearson correlation of approximately $r=0.31$ and RMSE of approximately $3.17$~kcal/mol against the experimental $\Delta G$ values. On the preserved antibody--antigen subset, only four matched cases remain, and the corresponding correlation is again weak. In addition, the log file \path{src/bridging/generatedData/PRODIGY/prodigy_failures.log} contains many execution failures, which made this branch operationally awkward as well as scientifically unconvincing.

For those reasons, PRODIGY was kept as an archived baseline branch rather than developed into the main thesis contribution. The underlying idea still made sense, but it did not provide a strong enough foundation for the final story. In particular, a correction-learning thesis built on a weak baseline would have been much harder to motivate than a direct comparison between static and MD-derived graph representations.

\subsection{MM/GBSA}

The MM/GBSA branch is preserved under \path{ARCHIVE/MMGBSA}, or on GitHub at \url{https://github.com/RohanDevereux/Bridging/tree/main/ARCHIVE/MMGBSA}. This branch includes AmberTools-based input preparation, sharded HPC launchers, prefetched estimate handling, and a runner that can execute MM/GBSA either from a single structure or from a supplied trajectory. The motivation here was again a correction-learning or comparison-baseline idea, but with a stronger physical flavour. MM/GBSA occupies an interesting middle ground in the literature: it is more physics-based than purely empirical scoring, but much cheaper than rigorous alchemical free-energy calculations \citep{Genheden2015MMGBSA}. That makes it a natural candidate baseline when one wants to ask whether machine learning can improve on an established approximate method.

The attraction of this path was obvious. If an endpoint free-energy approximation already captures some part of the binding signal, then learning a residual around that estimate could be more data efficient than learning affinity directly from raw structural descriptors. It also would have provided a clear bridge between classical computational chemistry and machine learning, which would have been thematically appealing for this project.

In practice, though, this branch never matured into a clean final result set comparable to the GraphVAE matrix reported in the main body of the thesis. The implementation is real and substantial, but the repo does not contain a final preserved metrics snapshot for MM/GBSA in the same directly auditable form as the final GraphVAE results. That matters. For the final thesis, I wanted the central empirical claims to be tied to artifacts that were still present, inspectable, and rerunnable from the repository state. MM/GBSA did not meet that bar by the time the project direction settled.

There is also a scientific reason why this branch became less attractive over time. As reviewed by Genheden and Ryde, MM/PBSA and MM/GBSA are useful but heavily approximation-dependent methods whose performance can vary strongly with the system, the solvent model, the entropy treatment, and implementation details \citep{Genheden2015MMGBSA}. That does not make them bad methods, but it does mean they are a difficult backbone for a thesis that also aims to compare machine-learning representations in a controlled way. In the end, MM/GBSA remained valuable as an idea and as preserved engineering work, but it was not the right centre of gravity for the final thesis.

\section{The C$\alpha$ contact-map branch}

The oldest major modelling branch in the repository is the contact-map pipeline preserved under \path{ARCHIVE/CA_contact_map}, or on GitHub at \url{https://github.com/RohanDevereux/Bridging/tree/main/ARCHIVE/CA_contact_map}. This branch converts selected interface residues into regular two-dimensional tensors built from C$\alpha$--C$\alpha$ contact and distance channels, then trains a convolutional variational autoencoder and downstream regression model on those tensors.

The motivation for this branch was straightforward. A contact map gives a fixed-size, image-like object that is much easier to process with standard convolutional architectures than a variable-size protein graph. That design choice is common in adjacent areas of structural machine learning, especially in protein--ligand binding work where voxelized or tensorized spatial representations have been used successfully, for example in Pafnucy \citep{Stepniewska2018Pafnucy}. For a project that was still trying to find a stable representation, this seemed like a sensible starting point: it imposed a strong geometric regularity on the data and made the learning problem easier to package.

The preserved code shows exactly what this branch was doing. Interface residues were selected from two chain groups using either frame-0 geometry or a more stable early-trajectory criterion in \path{ARCHIVE/CA_contact_map/src/bridging/featurization/interface.py}. The resulting residue subsets were converted into soft contact and clipped distance channels in \path{ARCHIVE/CA_contact_map/src/bridging/featurization/contact_map.py}. A set-based convolutional variational autoencoder was then trained on these tensors in \path{ARCHIVE/CA_contact_map/src/bridging/ml/model.py}, followed by downstream regression in \path{ARCHIVE/CA_contact_map/src/bridging/regression}.

This was not a nonsense design. In fact, it solved several engineering problems cleanly. Fixed tensor shape makes batching easy. The representation is compact. The interface selection logic is explicit. The model class is relatively small and interpretable. The problem is that this convenience comes at a scientific cost. By projecting a heterogeneous protein--protein interface onto a fixed C$\alpha$ contact tensor, the branch discarded a large amount of residue identity, side-chain, and graph-topological information. It also made the exact representation depend heavily on patch-selection and ordering choices. Once the graph pipeline was in place, the contact-map branch started to look like an over-compressed view of the same problem rather than the right abstraction for it.

So this branch was abandoned not because it was broken, but because it became hard to justify as the central representation once a residue-graph alternative existed. The graph formulation retained more chemically meaningful structure, handled variable-size complexes more naturally, and ended up being easier to connect to the MD-derived features that became the core of the final project.

\section{The PPB-inspired augmentation and spatiotemporal branch}

The PPB-inspired deep-learning branch is preserved under \path{ARCHIVE/PPB_Affinity_augmentation}, or on GitHub at \url{https://github.com/RohanDevereux/Bridging/tree/main/ARCHIVE/PPB_Affinity_augmentation}. This is probably the largest archived branch after the final pipeline itself. It contains three related modes: frame-level training, trajectory-level training, and an ensemble that combines predictions from several learned views. It also vendors a substantial external model stack under \path{ARCHIVE/PPB_Affinity_augmentation/src/BaselineModel}.

The motivation here was different from the PRODIGY and MM/GBSA branches. This branch was not mainly about correcting a baseline. It was about taking the idea of MD-informed learning more directly and more aggressively. The PPB-Affinity dataset paper explicitly argues for AI-based protein drug discovery on protein--protein complexes \citep{Liu2024PPBAffinity}, and more recent protein--ligand work has shown that spatio-temporal learning directly on MD-derived examples can be effective in at least some settings \citep{Libouban2025Spatiotemporal}. That made it tempting to try a more direct learning pipeline over frames and trajectories rather than reducing the MD to a small set of handcrafted graph features.

The archived implementation shows that this branch went well beyond a sketch. The frame mode in \path{ARCHIVE/PPB_Affinity_augmentation/src/bridging/ppb/frame_mode.py} trains a supervised structural model on sampled frame-level examples, with patch selection and augmentation. The trajectory mode in \path{ARCHIVE/PPB_Affinity_augmentation/src/bridging/ppb/trajectory_mode.py} uses a temporal model built on top of frame embeddings, including a bidirectional GRU over frame sequences. The ensemble mode in \path{ARCHIVE/PPB_Affinity_augmentation/src/bridging/ppb/ensemble_mode.py} fits a linear combiner across baseline, frame, and trajectory predictions. In other words, this branch represented a serious attempt at frame-aware and time-aware learning rather than a minimal proof of concept.

Why was it not carried through as the final thesis? The honest answer is that this was mainly a project-management pivot rather than a clean negative result. The branch was ambitious and interesting, but it also introduced several layers of complexity at once: a large vendored model stack, frame sampling decisions, temporal modelling decisions, augmentation policy, and evaluation logic distinct from the rest of the project. For a final-year thesis, that was a lot to own and explain rigorously. The later MD-to-graph pipeline was narrower, easier to audit, and easier to connect to concrete scientific questions such as ``does the full graph help more than an interface patch?'' and ``do these particular dynamic features improve held-out performance?'' Those questions fit the eventual scope of the thesis better.

So this branch should not be read as ``tried and failed.'' A more accurate reading is ``worth trying, partially implemented, but eventually superseded by a more focused and more explainable project direction.'' I think that distinction matters.

\section{Exploratory dataset curation utilities}

The final archived workstream is \path{ARCHIVE/Data_curation}, or on GitHub at \url{https://github.com/RohanDevereux/Bridging/tree/main/ARCHIVE/Data_curation}. These scripts cover earlier filtering decisions, partner-size prefetching, subset selection, diverse-sample construction, and older cleaning steps for intermediate datasets.

This material is less glamorous than the modelling branches, but it was still important. Before the final broad, pair-unique PPB-Affinity-derived dataset was fixed, there was a real question of which subset definitions were sensible and computationally feasible. Some branches focused on antibody--antigen cases, some on more diverse subsets, and some on excluding complexes whose partners were too large for a given representation or runtime budget. Those decisions directly affected what was possible to run, especially in the earlier stages of the project.

These scripts are archived simply because the final pipeline now has a stable dataset definition and no longer depends on those exploratory filters. They still matter for provenance, though, because they show that the final dataset was not chosen in a vacuum. It emerged from earlier attempts to balance biological scope, computational tractability, and representation design.

\section{What the archived branches contributed}

Although these branches did not survive into the final reported pipeline, they influenced the final thesis in concrete ways. The correction-learning branches clarified that fast structure-only or endpoint baselines were worth comparing against, but were not enough to carry the project on their own. The contact-map branch showed that enforcing a regular tensor shape made the problem easier to package but also removed too much useful structure. The PPB-inspired augmentation branch showed that direct frame- and trajectory-level learning was possible, but that it would be difficult to make the resulting thesis as auditable and focused as it needed to be. The data-curation utilities helped narrow the problem to a dataset definition that was broad enough to be interesting and structured enough to support repeated evaluation.

In that sense, the archived material is not just dead code. It is the record of how the final thesis question became sharper. The thesis eventually settled on a simpler but more defensible comparison: static versus MD-derived residue-graph features on a matched protein--protein affinity dataset, evaluated under repeated held-out resampling. That final formulation was reached by ruling out, narrowing, or deferring several other plausible directions, and the \path{ARCHIVE} directory is where that history now lives.
